% REPLACES: online_appendix.tex lines 1357-1609 ("\subsection{A.11. Price pooling when preparation can vanish} ..." to "... never as accepted rows.")
% This replaces the whole subsection with a short note. The subsection number,
% position, and label oa-a-pools stay, so A.10 and C.6 keep a valid target and no
% later section renumbers. Removed inside it: Proposition OA.3
% (label result-proposition-oa-3) and equations (OA.57)-(OA.64). The integrator
% decides whether later OA equation tags keep their numbers or close the gap.
% No \NUM placeholders.
\subsection{A.11. Price atoms under zero preparation
cost}\label{oa-a-pools}

With a positive low preparation cost, flows whose gross profit fell below that
cost could share a no-entry price, and the pooled posterior over that preimage
could support more than one price mapping at one order profile.
That construction has no counterpart under the declared cost law.
Gross profit is nonnegative at every belief on the reserve domain, and the
zero-cost realization prepares by the tie rule even where gross profit is zero.
Entry is therefore at least \(\rho\) at every price. In the reserve example
that the earlier construction used, gross profit at a candidate pooled
posterior \(\bar\mu\) was \(3\bar\mu>0\). A zero-cost buyer who saw that pooled
price would prepare, so the prescribed nonpreparation at the pooled price
fails at the first buyer check, before any pricing or trading check.

What remains is narrower. Where \(\Delta_T>0\), the increment argument in
Appendix A.7 of the paper ties the price to the raw posterior at fixed orders.
Distinct raw posteriors give distinct prices.
A price atom collects only flows with one common posterior, as on a Laplace
plateau \(x\ge1\) or \(x\le-1\) under full orders.
At such an atom the buyer's information set is still the whole preimage of the
price, and its posterior is the conditional probability over that preimage.
C.6 writes that formula as \eqref{eq:oa-oa-79} and validates each atom from its
full preimage. A candidate that let the buyer act on finer flow information
inside an atom would still fail, because that action is not measurable with
respect to the price.

Dow et al. (\citeproc{ref-DowGoldsteinGuembel2017}{2017}) discuss continuations
distinguished by their pooling regions following their Lemma 1.
In this model that distinction survives only through different order profiles,
and through degenerate reserves with no target-payoff spread, where the price
need not reveal beliefs. Continuation identity therefore keeps the whole-object
rule of \eqref{eq:oa-oa-53}: orders, price mapping with every atom, beliefs,
and preparation rule. No numerical exercise searches over alternative pooling
regions at fixed orders, and no registry quantity reports one.
